\input{../tex-inputs/latex-leseansicht-vorspann}

\section[Arthur und Olga Schnitzler an Gustav Schwarzkopf, 28. 4. 1904]{L04394 Arthur und Olga Schnitzler an Gustav Schwarzkopf, 28. 4. 1904}
\nopagebreak\mylabel{L04394v}
\rehead{ }\normalsize\beginnumbering\briefempfaengerindex{Schwarzkopf, Gustav@\textsc{Schwarzkopf, Gustav}!zzzSchnitzler, Olga@\emph{von Olga Schnitzler}!1904-05-281@{28. 4. 1904}|(be}\briefempfaengerindex{Schwarzkopf, Gustav@\textsc{Schwarzkopf, Gustav}!zzzSchnitzler, Arthur@\emph{von Arthur Schnitzler}!1904-05-281@{28. 4. 1904}|(be}
\toendnotes[C]{\smallbreak\pagebreak[2]}
\correspDesc{Versand  durch Arthur Schnitzler, Olga Schnitzler am 28. 4. 1904 in Ancona
\newline{}Zustellung  in Wien
\newline{}Erhalt  durch Gustav Schwarzkopf im Zeitraum [29. 5. 1904 – 2. 6. 1904?] in Wien}\toendnotes[C]{\smallbreak}
\Standort{DLA, A:Schnitzler, HS.1985.1.1897.}
\physDesc{Bildpostkarte, 173 Zeichen
\newline{}Handschrift Arthur Schnitzler: Bleistift, deutsche Kurrent
\newline{}Handschrift Olga Schnitzler: Bleistift, deutsche Kurrent
\newline{}Versand: Stempel: »\nobreak{}\oindex{I., Innere Stadt@\textbf{I., Innere Stadt}|pwk}Wien 1/1 1, bestellt\nobreak{}«.  }\toendnotes[C]{\smallbreak}\pstart{}{\pb}\textsc{Herrn Gustav Schwarzkopf}\pend{}\pstart{}\textsc{Vienna}\oindex{Wien@\textbf{Wien}|pw}\pend{}\pstart{}I. Tiefer Graben 23\oindex{Wien@\textbf{Wien}!I., Innere Stadt@\textbf{I., Innere Stadt}!Tiefer Graben 23@\textbf{Tiefer Graben 23}|pw}\pend{}\pstart{}\textsc{Austria.}\oindex{Österreich@\textbf{Österreich}|pw}\pend{}{\bigskip}
\pstart
           \centering{}{\pb}\textcolor{gray}{\textbf{Saluti da Ancona\oindex{Ancona@\textbf{Ancona}|pw}}}\pend
           
\pstart
           \centering{}\textcolor{gray}{\textbf{Panorama del Molo}}\pend
           \vspace{1em}
\pstart
           \raggedleft{}{\pb}28. 5. 904\pend
           \vspace{0.5em}
\pstart
           Herzliche Grüße, auf baldgs Wiederſehen\pend
           \pstart Ihr \spacefill\mbox{A.}\pend{}\selectlanguage{ngerman}\vspace{1em}
\pstart
           \noindent{}{[}hs. O. Schnitzler:{]} Wie freu ich mich auf’s \label{K_L04394-1v}\edtext{Heimfahren}{\lemma{\textnormal{\emph{Heimfahren}}}\Cendnote{\textnormal{Olga\pwindex{Schnitzler, Olga 17.\,1.\,1882 Wien – 13.\,1.\,1970 Lugano@\textsc{Schnitzler, Olga} (17.\,1.\,1882 Wien – 13.\,1.\,1970 Lugano)|pwk} und Arthur Schnitzler hatten Wien\oindex{Wien@\textbf{Wien}|pwk} am 30. 4. 1904 verlassen und waren den ganzen Mai durch Italien\oindex{Italien@\textbf{Italien}|pwk} gereist mit Stationen in Rom\oindex{Rom@\textbf{Rom}|pwk}, Neapel\oindex{Neapel@\textbf{Neapel}|pwk}, Palermo\oindex{Palermo@\textbf{Palermo}|pwk} und Taormina\oindex{Taormina@\textbf{Taormina}|pwk}. Am 30. 5. 1904 kehrten sie nach Wien\oindex{Wien@\textbf{Wien}|pwk} zurück.}}}\label{K_L04394-1}! \spacefill\mbox{Olga.}\pend
           
\pstart
           {[}hs. A. Schnitzler:{]} Ich glaube, ich auch. \spacefill\mbox{A.}\pend
           \selectlanguage{ngerman}\endnumbering\briefempfaengerindex{Schwarzkopf, Gustav@\textsc{Schwarzkopf, Gustav}!zzzSchnitzler, Olga@\emph{von Olga Schnitzler}!1904-05-281@{28. 4. 1904}|)be}\briefempfaengerindex{Schwarzkopf, Gustav@\textsc{Schwarzkopf, Gustav}!zzzSchnitzler, Arthur@\emph{von Arthur Schnitzler}!1904-05-281@{28. 4. 1904}|)be}\mylabel{L04394h}
\begin{anhang}
\end{anhang}\newcommand{\dateiname}{L04394}\newcommand{\titel}{Arthur und Olga Schnitzler an Gustav Schwarzkopf, 28. 4. 1904}\newcommand{\editorInnen}{Herausgegeben von Jahnke, SelmaMüller, Martin Anton}\input{../tex-inputs/latex-leseansicht-abspann}
